\section{光的衍射}
在过去初中的几何光学中，我们常说，\empx{光是沿直线传播}。这种说法没有太大的错误，但是更为深入的研究表明，光在通过小孔或狭缝等障碍物的边缘时会出现偏离直线传播的现象，例如使得单色光照射到开有圆孔的障碍屏上，当圆孔的直径较大时，光屏上的像将会呈现为一个均匀的圆形光斑，光斑也会随着圆孔的缩小而缩小，这时\uwave{光的直线传播}规律是成立的。但是如果将圆孔直径进一步缩小，使得圆孔直径和光的波长具有相近的数量级，此时如果进一步缩小圆孔直径，光斑就不再缩小了，光斑反而会变大，并且在其周围还出现一圈圈明暗相间的条纹。很明显该情形下，光的直线传播规律不再适用，光在该尺度上表现出了明显异于几何光线的波动特性。虽然尚不清楚原理，但姑且将光的这种特性，或者说传播方式，称为\uwave{光的衍射}。

光的干涉和光的衍射是两类不同的事物，前者是两束或多束相关光相遇时的一种现象，后者则是光本身具有的一种传播规律，光的衍射在一定意义上其实可以视为对光的直线传播规律在更小尺度上的修正。当然，干涉和衍射并非是割裂的，后面会看到两者实际具有密切联系。

\subsection{惠更斯--菲涅尔原理}
光学的历史上，正是衍射问题使光的波动说决定性战胜了光的微粒说，在法国科学院于1818年举办的征答辩论中，年仅30岁的菲涅尔创新的应用了惠更斯原理，圆满的解释了光在圆孔和狭缝等物体边缘的衍射现象，取得了辩论的胜利，从而使得光的波动学说得到最终的确认。

\uwave{惠更斯原理}是我们熟悉的，波前的每一点可以认为是产生球面次波的点波源，而之后任意时刻的波前可以视为这些次波的包络线。菲涅尔则进一步发展了惠更斯原理的想法，其进一步指出，\empx{空间中某点的振动状态是由前一时刻各次波在该点引发的相干叠加确定}。即，菲涅尔在惠更斯的“球面次波”的概念基础上进一步指出“球面次波会相互发生干涉”，这也就是所谓的\uwave{惠更斯--菲涅尔原理}，它解释了衍射线性的本质，\empx{光的衍射是波前球面次波的干涉结果}。

惠更斯--菲涅尔原理对应的数学形式是\uwave{菲涅尔衍射积分公式}，它的大致想法就是将无数的波前次波在场点上的作用通过积分叠加起来，就像过去我们在双缝干涉中做的那样。另外积分中还有一些相关的系数，称为\uwave{倾斜因子}，记为$K(\phi)$，它表示次波具有一定的方向性。换言之，虽然说次波是球面波，但是它对于球面并不是各向同性的，在原先的波动方向上次波显然具有最大的分量，在偏一些的方向上次波也就相对弱一些，最重要的是，次波并不会向后方传播。

由于数学上的复杂性，我们这里不给出菲涅尔衍射积分公式，一方面公式本身比较复杂，一方面即便有了公式在最简单情况下的定量计算都很繁琐，作为代价，我们后面在研究衍射现象时只能采取半定量手段，我们称为\uwave{半波带法}，它只能确定衍射图样的暗纹位置，而不能给出准确的光强分布函数，后者我们将作为结论不加证明的列出，并以其为依据绘制相应的图像。

\subsection{衍射的类型}
光的衍射物体通常可以归纳为下面两种类型，如\xref{fig:衍射的类型}所示
\begin{enumerate}
    \item 衍射障碍物与光源和光屏的距离均为有限远，称为\uwave{菲涅尔衍射}。
    \item 衍射障碍物与光源和光屏的距离均为无限远，称为\uwave{夫琅禾费衍射}。\footnote{关于“夫琅禾费衍射”中的“禾”的读音，牢记古诗：锄\udotcn{禾}日当午，汗滴\udotcn{禾}下土。}
\end{enumerate}
\begin{Figure}[衍射的类型]
    \begin{FigureSub}[菲涅尔衍射]
        \includegraphics[scale=0.75]{build/Chapter03E_04.fig.pdf}
    \end{FigureSub}\hspace{0.4cm}
    \begin{FigureSub}[夫琅禾费衍射]
        \includegraphics[scale=0.75]{build/Chapter03E_05.fig.pdf}
    \end{FigureSub}
\end{Figure}
这里我们关于\xref{fig:衍射的类型}做一些说明，首先比较让人困惑的是为何光线经过狭缝后，在狭缝的每个点上产生了三束不同方向的光线？这其实就是对惠更斯--菲涅尔原理的反应，即波前上的每个点都是新的次波源，而波前的各个次波源之间会相互发生衍射，这里绘制的不同方向的光线表示的就是作为球面波的次波源的波线，而各次波源沿同一方向的波线最终会汇聚于光屏。

这里夫琅禾费衍射中，我们要求衍射障碍物与光屏和光源的距离都是无穷远，因此，我们相应的使用了平行光源，其产生的衍射也是平行光，然而这就会导致了一个问题，平行光是无法成像的，为此我们在衍射障碍物和光屏间添加了一个凸透镜使光线会聚，从而在光屏上成像。

夫琅禾费衍射和菲涅尔衍射并不是孤立的，两者的判据是下面的\uwave{菲涅尔系数}
\begin{Equation}
    F=\frac{a^2}{L\lambda}\ll 1\qquad
    F=\frac{a^2}{L\lambda}\geq 1
\end{Equation}
前者将给出夫琅禾费衍射，后者将给出菲涅尔衍射，其中$a$表示圆孔半径或狭缝宽度，或比较一般的说是障碍物的线度，$\lambda$表示波长，$L$表示障碍物和光源以及光屏间的距离，因此这就很好理解了，当$F\ll 1$，如果$L$相对于障碍物线度$a$很大，那么$L$就可以视为无限远了，这就是夫琅禾费衍射，相反，当$F\geq 1$，就说明$L$相对障碍物线度是有限值，即菲涅尔衍射。

由于夫琅禾费衍射是相对比较简单，我们后面关于衍射的所有讨论都仅限于夫琅禾费衍射。

\subsection{单缝夫琅禾费衍射}
本小节我们来研究单缝的夫琅禾费衍射，首先我们不加证明的给出单缝衍射的光强分布函数。
\begin{BoxFormula}[单缝衍射的光强]
    单缝衍射在光屏角度相对中心$\phi$处产生的光强为
    \begin{Equation}&[]
        I=I_0\qty(\frac{\sin\alpha}{\alpha})^2
    \end{Equation}
    其中$\alpha$为代换变量
    \begin{Equation}&[代换变量]
        \alpha=\frac{\pi a}{\lambda}\sin\phi
    \end{Equation}
    其中，$a$为单缝宽度，$I_0$为单缝光强，$\lambda$为光的波长。
\end{BoxFormula}

根据\fancyref{fml:单缝衍射的光强}，我们可以绘制光强的函数曲线，如\xref{fig:单缝衍射的光强}所示。
\begin{Figure}[单缝衍射的光强]
    \includegraphics[scale=0.95]{build/Chapter03E_02.fig.pdf}
\end{Figure}
根据函数关系我们可以仿真出单缝衍射的效果，如\xref{fig:单缝衍射的光强仿真}所示。
\begin{Figure}[单缝衍射的光强仿真]
    \includegraphics[scale=0.95]{build/Chapter03E_01.fig.pdf}
\end{Figure}
现在让我们来看，如何运用我们知识体系范围内的方法，半定量的解释上述结果。

\begin{BoxFormula}[单缝衍射的暗纹分布]
    单缝衍射的暗纹分布如下，其中$\phi$为相对中心的角度
    \begin{Equation}&[]
        \phi=\pm(2m)\frac{\lambda}{2a}\qquad m\in\Z^{+}
    \end{Equation}
    其中，$a$为单缝宽度，$\lambda$为光的波长。
\end{BoxFormula}
\begin{Proof}
    我们使用的半定量分析方法称为\uwave{半波带法}，考虑\xref{fig:半波带法}。
    \begin{Figure}[半波带法]
        \includegraphics[scale=0.65]{build/Chapter03E_06.fig.pdf}
    \end{Figure}
    这里我们考虑衍射角为$\phi$的一组平行光束，它们经过透镜聚焦后会聚在光屏的点$P$上，透镜是具有等光程性的，这就是说，\empx{透镜物方焦面至像方焦面间的任意光路都是等光程的}，这很容易理解，从透镜两侧通过的光线走折线，几何路程较长，从透镜中心通过的光线走直线，几何路程较短，但是由于透镜中间厚两边薄的结构，从透镜中心通过的光线在高折射率的玻璃中通过的路程较多，这就补偿了光程，其在几何路程上“节约”的光程上“消耗”在玻璃中了。

    这里由于平行光是以$\phi$的角度通过透镜的，相应的可以认为光轴也偏转了$\phi$的角度，因此各束衍射光线在垂直光轴的$AH$面和光屏上的像方焦点$P$间的光程都是一致的，从而，它们间的光程差就完全来自$BH$部分了，由于这里$AH\perp BH$，因此$\angle BAH=\phi$，这样就有
    \begin{Equation}&[1]
        \delta=BH=AB\sin\phi=a\sin\phi
    \end{Equation}
    这里的$\delta$仅是$A$和$B$两点发出的衍射光间的光程差，而我们要计算的是$AB$上所有的点的衍射光之间的干涉结果，为此，我们可以将$\delta=BH$以$\lambda/2$为单位分为若干个\uwave{半波带}，很显然的是，在$AB$相邻的两个半波带上产生的衍射光线将会相互抵消，在\xref{fig:半波带法}中我们用红色和蓝色对半波带交替染色标识，这样，如果光程差刚好是$\lambda/2$的偶数倍，即有偶数个半波带
    \begin{Equation}&[2]
        \delta=\pm(2m)\frac{\lambda}{2}\qquad m\in\Z^{+}
    \end{Equation}
    此时，所有的半波带都相互抵消了，这就形成了暗纹，联立\xrefpeq{1}和\xrefpeq{2}即得
    \begin{Equation}&[3]
        a\sin\phi=\pm(2m)\frac{\lambda}{2}
    \end{Equation}
    由于这里在近轴成像，因此$\phi$很小，故可以作$\sin\phi=\phi$的近似
    \begin{Equation}&[4]
        a\phi=\pm(2m)\frac{\lambda}{2}
    \end{Equation}
    将$a$移项即得\xrefpeq{}的结论。
\end{Proof}

我们知道，暗纹的中间总是明纹，因此\fancyref{fml:单缝衍射的光强}可以推出明纹宽度。
\begin{BoxFormula}[单缝衍射的明纹宽度]
    单缝衍射的中央明纹角宽度如下
    \begin{Equation}
        \delt{\phi_0}=\frac{2\lambda}{a}
    \end{Equation}
    单缝衍射的次级明纹角宽度如下
    \begin{Equation}
        \delt{\phi}=\frac{\lambda}{a}
    \end{Equation}
    其中，$a$为单缝宽度，$\lambda$为光的波长。

    以上的结论表明，单缝衍射的中央明纹的宽度是次级明纹宽度的两倍。
\end{BoxFormula}
\fancyref{fml:单缝衍射的明纹宽度}的结论会让我们比较困惑，这里做一些解释
\begin{itemize}
    \item 首先，为何中央明纹具有不同于次级明纹的宽度？从数学上很容易看出原因在于衍射暗纹分布$\phi=\pm 2m(\lambda/2a)$并不适用于$m=0$的情形，即中央并不是暗纹，但为何中央就偏偏不是暗纹呢？关键在于，这里的$2m$描述的是半波带的个数，而不是先前在干涉中两束光的光程差相对半波长的倍数，后者在$m=0$意味着相等的光程差，并不特殊，前者在$m=0$则意味着零个半波带，这是极为特殊的，因为此时所有光线都具有相同的光程，因而根本不会发生任何的干涉相消，而是直接叠加在中央给出一个额外的明条纹。\nopagebreak
    \item 其次，为何中央明纹具有明显强于次级明纹的光强？如\xref{fig:单缝衍射的光强}所示，中央明纹的光强要显著高于哪怕是最近的次级明纹，从理论上可以证明，中央明纹集中了衍射几乎近$90\%$的能量。这种差异的根本原因，在于中央明纹和次级明纹的截然不同的形成方式
    \begin{enumerate}
        \item 衍射的中央明纹的成因是简单的，由于中央明纹处各个衍射光束的相位差是完全相同的，因此其不存在干涉相消的问题，这也就是说，中央明纹是整个单缝$AB$上的所有衍射光束的光强的叠加，通俗的说就是“劲往一处使”，所以中央明纹很亮。
        \item 衍射的次级明纹与暗纹的形成方式是类似的，都是干涉的结果，只不过说，后者是恰好有偶数个半波带的情形，前者则是对应奇数个半波带的情形，这种情况半波带在两两抵消后最终会多出一个半波带，这一部分半波带就构成了次级明纹的结果。很明显的是，随着衍射角的增大半波带的数量也会越来越多，因此最终抵消剩下的半波带在$AB$上的宽度也会越来越短，故衍射角越大，次级明纹的光强就越低。但即便在最好的情况下，次级明纹的光强也只可能来自$1/3$的$AB$段（即$3$个半波带的情形，因为$1$个半波代仍是中央明纹），所以，次级明纹的光强远不如中央明纹。
    \end{enumerate}
    \item 最后，还要说明的一点是，\empx{次级明纹的光强峰值其实并非在奇数个半波带时取到}。因为在奇数个半波带时虽然确实会多出一个半波带，但是如果衍射角再稍稍小一些，虽然此时剩下的半波带会不足一个，但是在$AB$上对应的半波带宽度也会大一些，两者的平衡之下，就有可能会出现，衍射角在稍小于奇数个半波长时，衍射的次级明纹反而取到极值。不过这种偏差并不是很大，这在\xref{fig:单缝衍射的光强}标识次级明纹极值的红色虚线中可以看出。
\end{itemize}
这里将次级明纹分布的定量结论列举如下，可以看出，半波长的奇数倍是相当好的近似。
\begin{BoxFormula}[单缝衍射的明纹分布]
    单缝衍射的次级明纹分布如下，其中$\phi$为相对中心的角度
    \begin{Equation}
        \phi=\pm(2m+1)\frac{\lambda}{2a}\qquad m=0.93,1.96,2.97,\cdots
    \end{Equation}
    在精度要求不高时可以近似为
    \begin{Equation}
        \phi=\pm(2m+1)\frac{\lambda}{2a}\qquad m\in\Z^{+}
    \end{Equation}
    其中，$a$为单缝宽度，$\lambda$为光的波长。
\end{BoxFormula}
这里还有一个值得注意的问题是，单缝衍射的中央明纹的角宽度是$\delt\phi_0=(2\lambda/a)$，如果光的波长$\lambda$相较单缝宽度$a$足够小，那么$\delt\phi_0\to 0$，这就表明衍射场基本集中在沿直线传播的原方向，在光屏上收缩为几何光学的像点，此时，光的衍射效应就可以忽略了，光的直线传播规律重新适用，由此可以看出，\empx{几何光学的实质就是波动光学在短波段$\lambda\to 0$时的极限形态}。

\subsection{圆孔夫琅禾费衍射}
圆孔夫琅禾费衍射的图样称为\uwave{艾里斑}，其明暗分布与单缝衍射的形态是类似的，由一组明暗相间的同心圆环围绕着中央的一个明亮的圆形亮斑（其占有约$84\%$的光能），如\xref{fig:艾里斑}所示。

然而，艾里斑的光强分布并不只是单缝衍射的光强分布在半径上的重复，艾里斑的光强分布函数是完全不同的，其中还会涉及到特殊函数论中的一阶贝塞尔函数$\mathrm{J}_1(x)$，陈列如下。
\begin{BoxFormula}[圆孔衍射的光强]
    圆孔衍射在光屏相对中心$\phi$处产生的光强为
    \begin{Equation}
        I=I_0\qty(\frac{2\mathrm{J}_1(\alpha)}{\alpha})^2
    \end{Equation}
    其中$\alpha$为代换变量
    \begin{Equation}
        \alpha=\frac{\pi a}{\lambda}\sin\theta
    \end{Equation}
    其中，$a$为圆孔直径，$I_0$为圆孔光强，$\lambda$为光的波长。
\end{BoxFormula}

\begin{Figure}[艾里斑]
    \includegraphics[scale=0.75]{build/Chapter03E_03.fig.pdf}
\end{Figure}
这里我们将圆孔衍射和单缝衍射的光强分布函数进行一下对比，如\xref{tab:圆孔衍射和单缝衍射的光强分布函数}所示，从中很容易看出，圆孔衍射的中央明纹的角宽度要比单缝衍射稍稍大一些，这也可以在以下公式中反映。
\begin{BoxFormula}[圆孔衍射的明纹宽度]
    圆孔衍射的中央明纹角宽度如下
    \begin{Equation}
        \delt\phi_0=1.22\frac{2\lambda}{a}
    \end{Equation}
    其中，$a$为圆孔直径，$\lambda$为光的波长。
\end{BoxFormula}

这里还有一个比较重要的问题是光学仪器的分辨本领，我们知道，许多光学仪器都可以视为圆孔，例如相机镜头或显微镜，由于衍射效应，当两个物点发出的平行光线通过透镜时，其在透镜的焦平面上形成的是两个艾里斑，而不是两个清晰的像点，这样导致的结果就是如果物点距离太近，两者衍射所形成的艾里斑将融为一体，此时，两者也就无法被光学仪器分辨了。\goodbreak

\begin{Table}[圆孔衍射和单缝衍射的光强分布函数]{|c|}
<圆孔衍射\\ \hlinelig>
\xcell<c>[1.5ex][0.0ex]{\includegraphics[scale=0.7]{build/Chapter03E_07a.fig.pdf}\hspace*{0.3cm}}\\
\hlinemid
单缝衍射\\ \hlinelig
\xcell<c>[1.5ex][0.0ex]{\includegraphics[scale=0.7]{build/Chapter03E_07b.fig.pdf}\hspace*{0.3cm}}\\
\end{Table}

那么，光学仪器的\uwave{最小分辨角}应该如何计算呢？我们有\uwave{瑞利判据}指出，如果一个艾里斑的中心恰好与另一个艾里斑的第一暗环的中心相互重合，此时恰能分辨两个艾里斑。很显然，这里的最小分辨角就是艾里斑\footnote{前面的解释有些不严密，艾里斑指的其实并非是整个圆孔衍射图案，艾里斑是特指其中的中央亮斑。}的半角宽$\delt\phi_0/2$，即第一暗环的半径（对应的角），我们记作$\fdd\theta$
\begin{Equation}
    \fdd\theta=1.22\frac{\lambda}{a}
\end{Equation}
因此，如果我们要增大光学仪器的最小分辨角$\fdd\theta$，就有两个方向。第一个方向是增大光学仪器的镜头口径，典型的例子是天文望远镜，现代天文台的光学天文望远镜的直径甚至可以达到米一级，目的就是为了增大对遥远天体的最小分辨角。第二个方向是减小光的波长，当然更一般的说，我们甚至都没有必要使用可见光乃至电磁波，典型的例子就是电子显微镜，电子束在德布罗意物质波的观点下相当于波长为$10^{-3}\si{nm}$的电子波，它的放大倍数可以达到数百万。